\documentclass[]{article}

\usepackage{amsmath}
\usepackage{multirow}
\usepackage{natbib}
\usepackage{graphicx} 


\begin{document}

The development of solid sorbents for industrial CO₂ capture is hindered by the conflicting requirements of strong chemical affinity for capture, low-energy thermal regeneration, and long-term structural durability. To identify materials that resolve these trade-offs, we present a high-throughput computational screening using the Materials Project database, systematically identifying 889 unique metal oxide-carbonate reaction pairs filtered for thermodynamic accessibility. Each candidate was evaluated against a comprehensive set of performance metrics, including Gibbs free energy to assess thermodynamic reversibility, volumetric expansion to predict mechanical integrity, and Tamman temperature to estimate sintering resistance. Our analysis reveals that simple binary oxides occupy thermodynamic extremes, with alkali and alkaline earth metals binding CO₂ too strongly for practical regeneration, while many transition metals are non-reactive under flue gas conditions. Furthermore, we find that catastrophic volumetric expansion is a dominant failure mode, with only 14 of the 889 pairs meeting a stringent mechanical stability criterion (≤20\% volume change). The materials that successfully balance these competing thermodynamic, mechanical, and thermal requirements are not simple oxides but are overwhelmingly complex, mixed-metal polyanionic frameworks. Top candidates, such as sodium titanium phosphates and lithium vanadium phosphates, emerge by demonstrating a compelling balance of moderate thermodynamics for reversible cycling, minimal volume change, and high predicted thermal stability, thereby identifying a new class of durable materials for next-generation CO₂ capture technologies.

\end{document}
                